\section{Introduction}
\label{sec:introduction}

Electroencephalography (EEG) records neural activity at the scalp with millisecond resolution, but the same electrodes also pick up the potentials generated by blinks and eye movements. These ocular artifacts are often larger than the neural activity of interest, spread across many channels, and overlap it in time and frequency \cite{croft2000ocular}. Rejecting every affected segment can discard a large fraction of a recording, while aggressive correction attenuates the event-related potentials and oscillations the recording was made to capture. Ocular-artifact removal is therefore a restoration problem: reduce the contribution of the eyes while leaving the neural signal intact. Classical regression solves it with an explicit physical model. When electrooculography (EOG) is recorded alongside EEG, the propagation of ocular potentials to each scalp channel can be estimated and the predicted contribution subtracted \cite{gratton1983ocular,semlitsch1986ocular}. Blind source separation and artifact subspace reconstruction remove the need for an EOG channel \cite{jung2000bss,chang2020asr}, and supervised networks trained on paired corpora learn far more flexible corrections than a fixed linear subtraction \cite{zhang2021eegdenoisenet,yu2022deepseparator}. Each family, however, fixes one side of the problem. Regression models the artifact pathway of one person but knows nothing about the structure of clean EEG. A trained network knows a great deal about clean EEG but averages the artifact pathway over its training population, and neither a subject label nor an embedding of the corrupted recording tells it how a particular person's eye activity actually reaches the electrodes.

The premise of this paper is that the subject-specific knowledge a denoiser needs is not an identity code. It is the operator that carries the subject: the EOG-to-EEG propagation matrix, which depends on anatomy, electrode placement, and session conditions, and which can be measured directly from about two minutes of simultaneous EEG and EOG \cite{gratton1983ocular,kobler2020sgeyesub}. We therefore split the denoiser into a part that is shared and a part that is measured. A conditional diffusion model, trained once on a population, represents the distribution of clean EEG \cite{ho2020ddpm,huang2025eegdfus}. A short calibration segment, recorded before the data to be cleaned, yields a ridge estimate of the user's propagation matrix. Empirical Bayes shrinkage pulls uncertain entries of this estimate toward the population value, and a reliability criterion decides whether the calibration can be used at all. At evaluation time the recorded EOG, passed through the calibrated matrix, becomes an explicit estimate of the ocular contribution that guides the diffusion model, and the same calibration supplies a per-channel feature vector that conditions the correction. Because the propagation estimate carries a known uncertainty, we sample it jointly with the diffusion trajectories, which turns the sampler into a source of predictive intervals whose width reflects both the model and the calibration.

Seen as a system, the method is a small adaptation loop rather than a fixed network: it measures the user's artifact pathway, checks the measurement, acts on it, and reports its remaining uncertainty. Each stage has a defined failure behavior. If the calibration segment is too short or unstable, the guide is withheld and restoration proceeds without subject-specific input, a fail-safe we adopted after measuring the alternative and finding that substituting a population matrix subtracts the wrong signal. Adapting to a new user requires no retraining and no gradient access, only a new calibration segment, and our measurements show that this dedicated calibration cannot be replaced by adapting the propagation estimate on-line during use. The predictive intervals close the loop by exposing, per channel and time point, how much the system trusts its own output.

We evaluate the method on multimodal recordings with 46 EEG channels and two bipolar EOG derivations, using paired episodes built from real EEG and real ocular waveforms for reconstruction error, and natural recordings for attenuation and preservation. On fifteen development participants, subject-calibrated guidance improved the relative root-mean-square error over an unguided control for every participant, and over population calibration for fourteen of fifteen. Eight further participants were held out until the models and the analysis were frozen and then evaluated once; the improvement reproduced at nearly identical magnitude. On natural recordings the calibrated system attenuated substantially more ocular signal than population calibration while preserving more of the recording. The predictive intervals reach nominal 80\% and 90\% coverage with a continuous ranked probability score slightly better than that of the uncalibrated sample distribution, so the calibration of uncertainty costs nothing in sharpness, and the sample-based intervals compare favorably with a deterministic ensemble. A supporting experiment on a second dataset with 259 training participants shows that the same participant-matched calibration information remains beneficial when delivered through a learned representation, and that the benefit neither grows nor shrinks with training-population size, which is consistent with our central premise: the useful subject information lives in the measured operator, not in the size of the population model.

This paper makes four contributions. First, a subject-calibrated EOG-guided diffusion method that combines closed-form propagation calibration, empirical Bayes shrinkage, and a reliability criterion with a fail-safe fallback, on top of a population diffusion model that is never retrained per user. Second, a controlled evaluation, with population, mismatched, unguided, and temporally shuffled comparisons, that locates the improvement in the user's own propagation estimate, together with a single-pass held-out evaluation that reproduces it. Third, natural-recording and deployment measurements, including the fallback comparison and an on-line adaptation comparison, that characterize how the system behaves when calibration is unreliable or unavailable. Fourth, calibration-aware predictive intervals that reach nominal coverage at no cost in probabilistic accuracy.
